\section{Conclusion}
\label{sec:conclusion}

We began by asking whether model quality can be predicted from weights alone---across architectures. Our work provides a nuanced answer: heavy-tailed self-regularization theory extends to Vision Transformer attention mechanisms, but cross-model comparison requires controlling for training provenance.

\subsection{Summary}

In this work, we conducted the first systematic measurement of heavy-tailed exponents on 53 Vision Transformer models from HuggingFace Model Hub. Our main contributions are:

\begin{enumerate}
    \item \textbf{Extension of HT-SR theory to Transformers.} Heavy-tailed exponents can be computed for ViT attention weight matrices ($\alpha \in [2.20, 18.92]$).

    \item \textbf{Identification of variance sources.} Training provenance---not architectural variation---dominates cross-model variance, with up to 12$\times$ variance reduction through family stratification (google/vit-*, n=6).

    \item \textbf{Practical guidance.} Weight-based model selection for diverse hubs must incorporate training origin metadata.
\end{enumerate}

\subsection{Future Directions}

\textbf{From untested alternative explanations.} Stratify models by training task to confirm whether task objective is the dominant variance factor.

\textbf{From unverified assumptions.} Re-evaluate a subset of models to verify HuggingFace accuracy metadata.

\textbf{From scope extensions.} Extend family analysis to ResNet, ConvNeXt, and Swin families.

\subsection{Closing Remarks}

As practitioners face the challenge of selecting among thousands of pretrained models, weight-based quality metrics offer a compute-efficient screening tool. However, these metrics are meaningful only with appropriate controls for model provenance. We hope this work encourages the research community to incorporate training metadata into cross-architecture analysis, moving toward truly universal weight-based model selection.
